\documentclass[10pt,a4paper]{amsart}
\usepackage{amsmath,amssymb,amsthm,graphicx}
\usepackage{verbatim}
\newenvironment{reference}{\comment}{\endcomment}
%\newenvironment{reference}{}{}
\newenvironment{slogan}{\comment}{\endcomment}
\newenvironment{history}{\comment}{\endcomment}

% For commutative diagrams we use Xy-pic
\usepackage[all]{xy}

% We use 2cell for 2-commutative diagrams.
\xyoption{2cell}
% Theorem environments.
%
\theoremstyle{plain}
\newtheorem{theorem}[subsection]{Theorem}
\newtheorem{proposition}[subsection]{Proposition}
\newtheorem{lemma}[subsection]{Lemma}

\theoremstyle{definition}
\newtheorem{definition}[subsection]{Definition}
\newtheorem{example}[subsection]{Example}
\newtheorem{exercise}[subsection]{Exercise}
\newtheorem{situation}[subsection]{Situation}

\theoremstyle{remark}
\newtheorem{remark}[subsection]{Remark}
\newtheorem{remarks}[subsection]{Remarks}

\numberwithin{equation}{subsection}

% Macros
%
\def\lim{\mathop{\mathrm{lim}}\nolimits}
\def\colim{\mathop{\mathrm{colim}}\nolimits}
\def\Spec{\mathop{\mathrm{Spec}}}
\def\Hom{\mathop{\mathrm{Hom}}\nolimits}
\def\Ext{\mathop{\mathrm{Ext}}\nolimits}
\def\SheafHom{\mathop{\mathcal{H}\!\mathit{om}}\nolimits}
\def\SheafExt{\mathop{\mathcal{E}\!\mathit{xt}}\nolimits}
\def\Sch{\mathit{Sch}}
\def\Mor{\mathop{\mathrm{Mor}}\nolimits}
\def\Ob{\mathop{\mathrm{Ob}}\nolimits}
\def\Sh{\mathop{\mathit{Sh}}\nolimits}
\def\NL{\mathop{N\!L}\nolimits}
\def\CH{\mathop{\mathrm{CH}}\nolimits}
\def\proetale{{pro\text{-}\acute{e}tale}}
\def\etale{{\acute{e}tale}}
\def\QCoh{\mathit{QCoh}}
\def\Ker{\mathop{\mathrm{Ker}}}
\def\Im{\mathop{\mathrm{Im}}}
\def\Coker{\mathop{\mathrm{Coker}}}
\def\Coim{\mathop{\mathrm{Coim}}}

% Boxtimes
%
\DeclareMathSymbol{\boxtimes}{\mathbin}{AMSa}{"02}

%
% Macros for moduli stacks/spaces
%
\def\QCohstack{\mathcal{QC}\!\mathit{oh}}
\def\Cohstack{\mathcal{C}\!\mathit{oh}}
\def\Spacesstack{\mathcal{S}\!\mathit{paces}}
\def\Quotfunctor{\mathrm{Quot}}
\def\Hilbfunctor{\mathrm{Hilb}}
\def\Curvesstack{\mathcal{C}\!\mathit{urves}}
\def\Polarizedstack{\mathcal{P}\!\mathit{olarized}}
\def\Complexesstack{\mathcal{C}\!\mathit{omplexes}}
% \Pic is the operator that assigns to X its picard group, usage \Pic(X)
% \Picardstack_{X/B} denotes the Picard stack of X over B
% \Picardfunctor_{X/B} denotes the Picard functor of X over B
\def\Pic{\mathop{\mathrm{Pic}}\nolimits}
\def\Picardstack{\mathcal{P}\!\mathit{ic}}
\def\Picardfunctor{\mathrm{Pic}}
\def\Deformationcategory{\mathcal{D}\!\mathit{ef}}
\usepackage{cleveref}
\crefname{originalreference}{original reference}{original references}
\Crefname{originalreference}{Original reference}{Original references}
\makeatletter
\newlabel{lemma-extend}{{01PF}{source}{}{sourcepointer-1}{}}
\newlabel{lemma-extend@cref}{{[originalreference][0][]01PF}{[1][0][]source}{}{}{}}
\makeatother
\crefname{originalreference}{original reference}{original references}
\Crefname{originalreference}{Original reference}{Original references}
\makeatletter
\newlabel{lemma-extend-trivial}{{01PE}{source}{}{sourcepointer-2}{}}
\newlabel{lemma-extend-trivial@cref}{{[originalreference][0][]01PE}{[1][0][]source}{}{}{}}
\makeatother
\crefname{originalreference}{original reference}{original references}
\Crefname{originalreference}{Original reference}{Original references}
\makeatletter
\newlabel{lemma-finite-presentation-module}{{01PC}{source}{}{sourcepointer-3}{}}
\newlabel{lemma-finite-presentation-module@cref}{{[originalreference][0][]01PC}{[1][0][]source}{}{}{}}
\makeatother
\crefname{originallemma}{original lemma}{original lemmas}
\Crefname{originallemma}{Original lemma}{Original lemmas}
\makeatletter
\newlabel{modules-lemma-kernel-surjection-finite-free-onto-finite-presentation}{{01BP}{source}{}{sourcepointer-4}{}}
\newlabel{modules-lemma-kernel-surjection-finite-free-onto-finite-presentation@cref}{{[originallemma][0][]01BP}{[1][0][]source}{}{}{}}
\makeatother
\title{Extending a finitely presented module with a map to a quasi-coherent module from a quasi-compact open}
\author{The Stacks Project Authors}
\date{stacks-project:properties.tex:local-snapshot-20261012; GNU FDL 1.2 or later}
\begin{document}
\maketitle
\noindent\textit{Editorial note.} All delivered mathematical text is the author's own, copied from the pinned Stacks Project chapter file properties.tex; no statement or proof step was rewritten, reordered or removed.\par\medskip
\noindent\textit{Editorial note.} The counted proof is the author's complete argument: the reduction to one affine piece outside U and then to an affine base by glueing, the production of a quasi-coherent extension with its shrinking to a finite type submodule, and the presentation by a finite free module with the control of the kernel giving finite presentation.\par\medskip
\noindent\textit{Editorial note.} The label lemma-extend-finite-presentation names the counted statement, and the proof environment immediately following it in the source is delivered verbatim as its proof.\par\medskip
\noindent\textit{Editorial note.} The counted statement is carried by the author's own lemma environment in the chapter source.\par\medskip
\noindent\textit{Editorial note.} Pre-open gate: python operations/agent-authoring/tools/precard\_gate.py --source properties.tex --label lemma-extend-finite-presentation --start 3140 --end 3226 --local-id 28033 returned RESULT=OK (L28033 free, occupancy FREE, public interval clear). Candidate taken from sources/stacks-properties/candidates-halfA.json.\par\medskip
\noindent\textit{Editorial note.} The statement and proof are copied from the pinned chapter file; the cited inputs are kept as external pointers with their canonical Stacks tags.\par\medskip
\noindent\textit{Allowed standard prerequisites (editorial).}\begin{itemize}
\item Extension of a pair by a quasi-coherent sheaf on an affine scheme. Same chapter, the trivial extension lemma producing a quasi-coherent sheaf with a map to the given one restricting to the given pair over the open. Supplies the sheaf H and the map psi on the affine case.
\item Extension of a finite type submodule from a quasi-compact open. Same chapter, the lemma producing a finite type quasi-coherent submodule restricting to a given quasi-coherent subsheaf over the open. Supplies the finite type submodules H' and K' used in the two reduction steps.
\item Description of quasi-coherent modules on an affine scheme and the kernel of a surjection. Same chapter, the lemma identifying a finite type quasi-coherent module on an affine scheme with the module attached to a finitely generated module, and Modules, the lemma on the kernel of a surjection from a finite free module onto a finitely presented module. Supply the surjection from a finite free module onto H and the finite type property of the kernel restricted to U.
\end{itemize}\medskip
\ifdefined\hypertarget\hypertarget{sourcepointer-1}{}\fi
\noindent\textit{Source reference.} Original reference 01PF: sources/stacks-properties/pinned-original-source.tex, source lines 3064--3064. Prerequisite pointer invoked at proof line(s) 3157; it names a standard fact of the same chapter that the counted argument uses, and every step of the counted proof built on it is retained verbatim in the delivered body.\par\medskip
\ifdefined\hypertarget\hypertarget{sourcepointer-2}{}\fi
\noindent\textit{Source reference.} Original reference 01PE: sources/stacks-properties/pinned-original-source.tex, source lines 3011--3011. Prerequisite pointer invoked at proof line(s) 3185; it names a standard fact of the same chapter that the counted argument uses, and every step of the counted proof built on it is retained verbatim in the delivered body.\par\medskip
\ifdefined\hypertarget\hypertarget{sourcepointer-3}{}\fi
\noindent\textit{Source reference.} Original reference 01PC: sources/stacks-properties/pinned-original-source.tex, source lines 2330--2330. Prerequisite pointer invoked at proof line(s) 3196; it names a standard fact of the same chapter that the counted argument uses, and every step of the counted proof built on it is retained verbatim in the delivered body.\par\medskip
\ifdefined\hypertarget\hypertarget{sourcepointer-4}{}\fi
\noindent\textit{Source reference.} Original lemma 01BP: sources/stacks-chow/tags.txt, source lines 1734--1734. Prerequisite pointer invoked at proof line(s) 3208; it names a standard fact of a sibling chapter that the counted argument uses, and every step of the counted proof built on it is retained verbatim in the delivered body.\par\medskip
\begin{lemma}
\label{lemma-extend-finite-presentation}
Let $X$ be a quasi-compact and quasi-separated scheme.
Let $\mathcal{F}$ be a quasi-coherent $\mathcal{O}_X$-module.
Let $U \subset X$ be a quasi-compact open.
Let $\mathcal{G}$ be an $\mathcal{O}_U$-module which is of finite presentation.
Let $\varphi : \mathcal{G} \to \mathcal{F}|_U$ be a morphism of
$\mathcal{O}_U$-modules.
Then there exists an $\mathcal{O}_X$-module
$\mathcal{G}'$ of finite presentation, and a morphism
of $\mathcal{O}_X$-modules $\varphi' : \mathcal{G}' \to \mathcal{F}$
such that $\mathcal{G}'|_U = \mathcal{G}$ and such that
$\varphi'|_U = \varphi$.
\end{lemma}

\begin{proof}
The beginning of the proof is a repeat of the beginning of the
proof of Lemma \ref{lemma-extend}. We write it out carefully anyway.

\medskip\noindent
Let $n$ be the minimal number of affine opens $U_i \subset X$,
$i = 1, \ldots , n$ such that $X = U \cup \bigcup U_i$.
(Here we use that $X$ is quasi-compact.) Suppose
we can prove the lemma for the case $n = 1$. Then we can successively
extend the pair $(\mathcal{G}, \varphi)$
to a pair $(\mathcal{G}_1, \varphi_1)$ over $U \cup U_1$
to a pair $(\mathcal{G}_2, \varphi_2)$ over $U \cup U_1 \cup U_2$
to a pair $(\mathcal{G}_3, \varphi_3)$ over $U \cup U_1 \cup U_2 \cup U_3$,
and so on.
Thus we reduce to the case $n = 1$.

\medskip\noindent
Thus we may assume that $X = U \cup V$ with $V$ affine.
Since $X$ is quasi-separated and $U$ quasi-compact,
we see that $U \cap V \subset V$ is quasi-compact.
Suppose we prove the lemma for the system
$(V, U \cap V, \mathcal{F}|_V, \mathcal{G}|_{U \cap V}, \varphi|_{U \cap V})$
thereby producing $(\mathcal{G}', \varphi')$ over $V$.
Then we can glue $\mathcal{G}'$ over $V$ to the given sheaf $\mathcal{G}$
over $U$ along the common value over $U \cap V$, and similarly we can glue
the map $\varphi'$ to the map $\varphi$ along the common value over
$U \cap V$. Thus we reduce to the case where $X$ is affine.

\medskip\noindent
Assume $X = \Spec(R)$.
By Lemma \ref{lemma-extend-trivial} above we may
find a quasi-coherent sheaf $\mathcal{H}$ with
a map $\psi : \mathcal{H} \to \mathcal{F}$ over $X$
which restricts to $\mathcal{G}$ and $\varphi$ over $U$.
By Lemma \ref{lemma-extend} we can find a finite type
quasi-coherent $\mathcal{O}_X$-submodule
$\mathcal{H}' \subset \mathcal{H}$
such that $\mathcal{H}'|_U = \mathcal{G}$. Thus after
replacing $\mathcal{H}$ by $\mathcal{H}'$
and $\psi$ by the restriction of $\psi$ to $\mathcal{H}'$
we may assume that $\mathcal{H}$ is of finite type.
By Lemma \ref{lemma-finite-presentation-module}
we conclude that $\mathcal{H} = \widetilde{N}$ with
$N$ a finitely generated $R$-module. Hence there exists a surjection
as in the following short exact sequence of
quasi-coherent $\mathcal{O}_X$-modules
$$
0 \to \mathcal{K} \to \mathcal{O}_X^{\oplus n} \to \mathcal{H} \to 0
$$
where $\mathcal{K}$ is defined as the kernel.
Since $\mathcal{G}$ is of finite presentation and
$\mathcal{H}|_U = \mathcal{G}$ by
Modules, Lemma
\ref{modules-lemma-kernel-surjection-finite-free-onto-finite-presentation}
the restriction $\mathcal{K}|_U$ is
an $\mathcal{O}_U$-module of finite type. Hence by Lemma \ref{lemma-extend}
again we see that there exists a finite type quasi-coherent
$\mathcal{O}_X$-submodule $\mathcal{K}' \subset \mathcal{K}$ such
that $\mathcal{K}'|_U = \mathcal{K}|_U$. The solution to the problem
posed in the lemma is to set
$$
\mathcal{G}' = \mathcal{O}_X^{\oplus n}/\mathcal{K}'
$$
which is clearly of finite presentation and restricts to give $\mathcal{G}$
on $U$ with $\varphi'$ equal to the composition
$$
\mathcal{G}' = \mathcal{O}_X^{\oplus n}/\mathcal{K}'
\to \mathcal{O}_X^{\oplus n}/\mathcal{K} = \mathcal{H} \xrightarrow{\psi}
\mathcal{F}.
$$
This finishes the proof of the lemma.
\end{proof}

\end{document}
